\documentclass[a4paper,12pt]{article}

% Storyline template by Jan van Gemert

% use URL clickable link
\usepackage{hyperref}

% for space betwen enumerate/itemize
\usepackage{enumitem}

% using color boxex
\usepackage{xcolor}

\title{The Storyline }
\author{Jan van Gemert \\ Delft University of Technology \\
	Full Research Guidelines are: \\ \url{https://jvgemert.github.io/ResearchGuidelinesInDL.pdf} }
\date{}
%\date{\vspace{3cm}  %\includegraphics[width=0.6\linewidth]{vraagroeptekenScheiding} \clearpage }

\begin{document}


\clearpage

Empty storyline template (replace with paper title).  \\
Please follow: \url{https://jvgemert.github.io/storyline.pdf}. \\
Each term that is used needs to be introduced/motivated. As a tool: Each \colorbox{cyan!20}{first-time} use of a term can be colored; a \underline{re-used term} underlined. \\
Try to fit to a single page; remove unused words; grammar is optional; change margins or font size if really needed. {\scriptsize OK; if 2 pages are needed; so be it, but rather not :).} \\



\begin{enumerate}[left=0pt,nosep]
	\item \textbf{Why interesting?}
	\begin{enumerate}[left=-10pt,label=\alph*,nosep]
		\item ...
		\item ...
	\end{enumerate}
	\item \textbf{How done now?}
	\begin{enumerate}[left=-10pt,label=\alph*,nosep]
		\item ...
		\item ...
	\end{enumerate}
	\item \textbf{What is missing, and So What?}
	\begin{enumerate}[left=-10pt,label=\alph*,nosep]
		\item ...
		\item ...
	\end{enumerate}
	\item \textbf{Proposed approach (P). }
	\begin{enumerate}[left=-10pt,label=\alph*,nosep]
		\item ...
		\item ...
	\end{enumerate}
	\item \textbf{Experimental questions.}
	\begin{enumerate}[left=-10pt,label=\alph*,nosep]
        \item ... \\
        c1: .\\
        c2: .\\
        c3: .\\
        c4: .
        \item ... \\
        u1: .\\
        u2: .\\
        u3: .\\
        u4: .
	\end{enumerate}
\end{enumerate}



\end{document}
